%!TEX program = uplatex
\RequirePackage{plautopatch}
\documentclass[uplatex,dvipdfmx]{jlreq}
\usepackage{amsmath,amssymb}

\pagestyle{empty}

\begin{document}

%======================================================================
% 表紙
%======================================================================
\begin{center}
  \vspace*{20pt}
  {\Huge\bfseries Web 幾何学}

  \vspace{30pt}

  {\Large\bfseries リーマン面・微分幾何・計算図表・ツイスター理論}

  \vspace{50pt}

  {\large 前文}

  \vspace{10pt}

  \parbox{0.85\textwidth}{
    \large\bfseries
    Web 幾何学の歴史的発展と現代的展開を概観する。
    Translation 面の幾何学から始まりカルタン、ブラシュケの微分幾何を経て、
    計算図表・ツイスター理論との関係まで幅広く論じる。
  }
\end{center}

\vfill

\newpage

%======================================================================
% 講演１：中居功
%======================================================================
\noindent
{\large
\begin{enumerate}
  \setlength{\itemsep}{2pt}
  \item なぜ２重 Translation 面からリーマン面がでてくるか？\\
        \hspace*{\fill}——リー、ポアンカレ、ヴィルティンガー、ダルブーの仕事
  \item Web の微分幾何の初歩：カルタン、ブラシュケ、チャーンから現在まで
  \item 幾何光学における Wave Front の Web 幾何学の試み
\end{enumerate}
}

\bigskip
\hspace*{\fill}{\large\bfseries 中居 功（北海道大学・数学教室）}

\bigskip

Web 幾何学とは１９２０年代に Blaschke により名づけられた、
（一時はハチの巣幾何学などと呼び誤解をまねいたこともあるが）
現代の用語によれば、いくつかの葉層の重ね合わせの幾何構造の研究の総称である。
なぜそのような物を考えるのだろうか？
素朴には、波は幾重にも重なっているとき豊かな構造をもつし、また微分方程式
$f(x,y,y')=0$ の解曲線は葉層の重ね合わせをなすからといえる。
が、意外にも最初の深い結果はリーマン面のヤコビアンの中のテータ因子の研究からでてきた。
それから１００年以上かかってゆっくりと着実に発展してきた。
Web 幾何学は Blaschke 以前は異なる興味のもとに様々な方法で研究されてきた。
この講義ではその中の大きな二つの歴史的トピックの紹介をする。
代数幾何以前のリーマン面の幾何学ではじまり、現代へととぶ。
また最後に、幾何光学における Wave Front のなす Web 構造がきわめて自然に
Web 幾何学の枠組みにあてはまることをいくつかの例をとって説明する。

\bigskip

\noindent\textbf{第一回目}では、
１８８０年ころの Lie の Translation 曲面の研究での大発見、
それに続く Poincaré、Wirtinger、Darboux の仕事を紹介する。
$n$ 次元ユークリッド空間の Translation 曲面とは、$n$ 本の空間曲線の点のベクトルとしての
和全体の集合 $C_1 + \cdots + C_n$ のことをいう。
この曲面は例えば古くは Monge によって詳しくしらべられている。

リーマン面 $C$ のヤコビアンの中に棲むテータ因子 $\Theta$ はもともと
標準積分曲線による葉層構造を備えているが、まさにそれが Translation 曲面の最も重要な例である。
少し詳しくいうと $\Theta$ は $C$ の $g-1$ 個の点の形式和
$D = x_1 + \cdots + x_{g-1}$ によってパラメトライズされるが、
これがベクトルの和による Translation 構造をあらわしている。
アーベルの定理は、$D + E$ が $C$ の正則１形式の零点集合（標準因子）であるとき
零ベクトルに恒等的に等しいことをいうが、これはテータ因子 $\Theta$ が $D$ と $E$ による
ふたつの Translation 構造を持つことを言っている。
このような曲面を２重 Translation 面という。
Lie の発見は、全ての局所的な２重 Translation 面はあるテータ因子 $\Theta$ の一部分であり、
また局所的２重 Translation 構造は曲面に対してあるとすれば唯一であることである。
これらの美しい結果も、Lie の証明が偏微分方程式を用い、また不透明であったこともあり、
すぐには理解されなかったが、
１０年程後に Poincaré、さらに後に Darboux によってよりよい証明がつけられた。
上の表現は彼等によるところが大きい。
（現在 Torelli の名で呼ばれているリーマン面のヤコビアンの定理もここでの議論から
すぐにわかるが、１８８０年には彼はまだ現われていない。）
とくに Darboux による証明は、現在も指針をあたえている。
ここにはいくつものすばらしいアイディアが輝いている。

\bigskip

\noindent\textbf{第二回目}では、
Cartan、Blaschke--Bol による Web の微分幾何の初歩と線形化の問題、
それにまつわるトピックを紹介する。
これは Web 幾何学のもう一つの生い立ちを知る上でかかせない。
Web の微分幾何はすでに Cartan の一変数の一階常微分方程式 $y' = f(x,y)$ の
直積形の変数変換 $(x,y) \to (\phi(x), \psi(y))$ による分類問題の研究の中にみられる。
この変数変換は解曲線による葉層のほかに座標関数による葉層を保つ。
Cartan は $xy$-平面上にアフィン接続を定義し、それが方程式の不変量であることを示した。

一般に、平面に $d$ 個の葉層があるとき、それを $d$-Web とよぶ。
それら葉層をなす曲線すべてが満たす２階常微分方程式 $y'' = f(x,y,y')$ は容易に書き下せるが、
$d$-Web の分類はその２階常微分方程式の分類とも考えられる。
ほぼ同じ１９世紀末に Tresse、Cartan の２階常微分方程式 $y'' = f(x,y,y')$ の研究があるが、
それによって $y'' = 0$ と同値になるための条件が知られている。
これは Web の線形化、つまり座標変形で Web の曲線すべてを直線にできるための条件にほかならない。
線形化できれば各曲線の双対をとることで双対射影空間の $d$ 個の曲線片ができるが、
Darboux の結果によれば、ある条件でそれが一本の代数曲線に拡張し（代数化）リーマン面ができあがる。
現在も一般次元での線形化は Web 幾何学の最も重要な問題であり未解決である。
Bol は平面の $d$-Web でアフィン接続が平坦なものを分類し、そのなかで
唯一、線形化できない例外 Bol Web とよばれる $5$-Web を発見し、
それの５つの葉層の定義関数の満たす関係式（Abel 方程式）として
Rogers dilogarithm のいまではよく知られた５項関係式が現われることを示した。
また例外 Bol Web は MacPherson、Damiano によって高次元化されている。

もう一つの展開として、アフィン接続を不変に保つ Web の変形の問題がある。
$1$-形式の線形族のなかで可積分な $d$ 個の $1$-形式のなす $d$-Web はすべて
共通のアフィン接続をもつが、そのような性質を持つものは線形族からくるものに限ることが最近証明できた。
$1$-形式の射影線形空間のなかで可積分条件をみたす Veronese 曲線を Veronese Web という。
Gel'fand--Zakharevich によれば Veronese Web と奇数次の bi-Hamiltonian system は
１対１関係にあることが知られている。

\bigskip

\noindent\textbf{第三回目}では、
幾何光学における Wave front（波面）の時間による Propagation の Web 幾何学を考える。
$\mathbf{R}^n$ 上の Wave front の Propagation は射影余接空間の $n$ 次元曲面 $S$ と
その上の時間関数による余次元１葉層で与えられる。
$S$ 上で時刻 $= t$ の超局面の底空間への射影が時刻 $t$ の波面となる。
例えば、ひだを寄せながら砂浜に寄せる波を思いうかべればわかるように、波面は互いに重なりあっている。
いいかえれば、$S$ 上の時刻による葉層は底空間に射影したとき、様々な Web 構造をなしている。
この Web の特異点集合が Caustics である。

$S$ での時刻による葉層の Bott 接続は自然な Chern 接続と呼ばれる $S$ の葉層を保アフィン接続に拡張する。
これを底空間に射影すると、そこでの平均接続の曲率形式、
つまり平均曲率が Blaschke 曲率となる。

この Blaschke 曲率形式が面白い！
それは波の位相を観察している空間上で微分幾何的にとらえているのだが、
その Caustics に沿って Web の特異性に応じて極をもつのである。
また Web の moduli 空間と、いくつかの例での観察ではあるが、１対１対応にある。
これらの枠組みはまだ、作ったばかりで中味はそう多くないが、
とても美しいと思うのは私だけだろうか？

\bigskip

\begin{center}
  \textsc{References}
\end{center}

\begin{enumerate}
  \item W.~Blaschke and G.~Bol,
        \textit{Geometrie der Gewebe},
        Grundlehren der Mathematischen Wissenschaften,
        Springer, 1938.
  \item I.~Nakai,
        Web 幾何学：ミクロコスモス,
        竜谷大学科学技術共同研究センター, 講究録, 1996, 29--98.
\end{enumerate}

\newpage

%======================================================================
% 講演２：佐藤肇
%======================================================================
\noindent
{\large\bfseries 計算図表、Web、微分方程式の標準化とツイスター理論}

\medskip
\hspace*{\fill}{\large\bfseries 佐藤 肇（名古屋大学・多元数理）}

\bigskip\bigskip

最近のデジタル型コンピューターの発展により影がうすくなったが、
もう一度アナログ型コンピューターを考えてみよう。
もう中年にしか記憶にないかもしれないが、計算尺を知っている人も多いであろう。
計算図表は計算尺をやや複雑にしたものであり、Web 理論の基となった。

\bigskip

\noindent\textbf{問題：}$(x,y,z)$ 空間での関数 $F(x,y,z)$ が与えられた時、
$F(x,y,z) = 0$ を満たす組 $(x,y,z)$ を求めよ。
あるいは $(x_0, y_0)$ に対して $F(x_0, y_0, z_0) = 0$ を満たす $z_0$ を探せといってもよい。

\bigskip

\noindent\textbf{例１}\quad $F(x,y,z) = z^2 + yz + x$

この時は、積が $x$、和が $-y$ となる２つの数 $z_1, z_2$ を求めよという問題となる。

\medskip

\noindent\textbf{例２}\quad $F(x,y,z) = x^2 + y^2 + z$

\bigskip

\noindent\textbf{解法１}\quad $F(x,y,\varphi(x,y)) = 0$ となる陰関数 $\varphi(x,y)$ を求めればよい。

\medskip

より図形的にこのような解を求める方法は次のようになるであろう。

\medskip

\noindent\textbf{解法２}\quad $F(x,y,z) = 0$ は空間の一つの曲面を定める。
その曲面の $(x_0, y_0)$ における高さが求める $z_0$ の値となる。

\medskip

\noindent\textbf{解法３}\quad さらに、$(x,y)$ 平面上で解を知るには、
曲面 $F(x,y,z) = 0$ の $z\,(= \varphi(x,y)) = \mathrm{const}$ という等高線を細かく引いて、
$x = x_0$、$y = y_0$ という直線と交わる等高線の示している高さ $z_0$ をみればよい。
この方法は、$F(x,y,z) = 0$ という曲面を与えると、
\begin{itemize}
  \setlength{\itemsep}{2pt}
  \item $x = \mathrm{const}$ という直線による一組の族、
  \item $y = \mathrm{const}$ という直線による一組の族、
  \item $z = \mathrm{const}$ という曲線による一組の等高線の族
\end{itemize}
が与えられているとみなしている。

これが共点図表で、いいかえれば、平面の $3$-Web でもある。

\bigskip

例１の場合を考えると、$z = \mathrm{const}$ という曲線もすべて直線になり、
３組の直線族による $3$-Web となる。
例２はそうではないが、
平面の $x \mapsto x^2$、$y \mapsto y^2$ という座標変換により直線族となる。

\medskip

このように、曲線の族による平面の $3$-Web を座標変換により直線族による $3$-Web と
みなすことができるかというのが、
共線図表（＝計算図表（狭い意味で））化可能性という計算図表学の基本問題であり、
また Web 理論・微分方程式論の中心問題でもある。
この解答は、ツイスター理論と密接な関係があることがわかる。

\end{document}
